\begin{textsample}{NEG Dim 3 – pa – Score -3.00 – pa/pa\_em\_31.txt}  \label{ex:f3_neg_005}
Quadro 8 : Correlação entre as categorias de área de Linguagens e suas tecnologias e o campo de saberes e práticas do ensino de Educação Física Categoria de área : Vida Pessoal Campo de saberes e práticas da Educação Física : Nessa categoria, os ( as ) professores ( as ) de Educação Física, de modo interdisciplinar e contextualizado, devem reunir esforços para desenvolver \textbf{competências} e habilidades específicas para esta etapa de ensino, previstas pela BNCC, que fomentem a construção de suas subjetividades, emoções, interpretações e identidades corporais.
Como exemplo, a cultura corporal pode ser trabalhada resgatando -se as experiências corporais de vida dos alunos, articulando, ampliando e aprofundando as diversas manifestações culturais e corporais que estabeleçam relações diretas e indiretas com o seu cotidiano, na busca pelo desenvolvimento do protagonismo e autonomia para com tais manifestações culturais, sejam elas ligadas ao jogo, dança, folclore, esporte, lutas, ginástica e outras, articuladas a macrotemáticas que busquem o debate sobre o referido objeto de conhecimento, selecionando em cada uma das etapas a melhor maneira de articulá -la com as demais áreas e/ou Campos de Saberes e Práticas.
Categoria de área : Práticas de estudo e pesquisa Campo de saberes e práticas da Educação Física : Esta categoria propicia a aprendizagem colaborativa, integra teoria e prática ao levar o aluno à condição de pesquisador e construtor do próprio conhecimento.
Para atuar com esta categoria, os ( as ) professores ( as ) de Educação Física precisam planejar projetos ou práticas curriculares que levem o aluno a pesquisar, analisar e pôr em prática os conhecimentos adquiridos.
Por exemplo, pode -se trabalhar as práticas de estudos e pesquisas ligadas à cultura corporal por meio do levantamento de dados estatísticos sobre diversas temáticas ligadas à área, aprofundamentos teóricos e debates que possibilitem a ampliação e percepção do jovem sobre os diversos discursos presentes nas práticas corporais.
Essa categoria possui uma forte conexão com outras áreas do conhecimento, como as Ciências Humanas, as Ciências da Natureza e a Matemática, uma vez que a pesquisa em Educação Física se dá na apropriação das representações do corpo nas diferentes formas do saber.
Outrossim, as demais áreas anteriormente citadas também devem se utilizar das práticas corporais, que muito atraem os jovens nesta etapa de ensino, como ferramenta para integrar os conhecimentos historicamente produzidos nas diferentes áreas do saber, com a consolidação dos elementos da pesquisa no cotidiano de vida dos alunos.
Categoria de área : Cultural artístico e literário Campo de saberes e práticas da Educação Física : Nessa categoria, os ( as ) professores ( as ) de Educação Física devem articular e proporcionar ao aluno o reconhecimento de sua identidade cultural.
Nesse sentido, deve -se ampliar a participação/reflexão nas atividades culturais, o senso crítico e a interação do jovem ao mundo artístico-cultural em diferentes contextos, tais como : teatros, cinemas, espaços poliesportivos e de movimentos corporais, artísticos e sociais que estejam relacionados com as variadas nuances em que se apresenta a dimensão cultural na atualidade ( visual, sonora, corporal, dramaturgia e midiáticas ), de modo a promover a fruição, produção, transformação, questionamento e a crítica das relações entre as representações corporais e expressivas, simbólicas e audiovisuais de posições subjetivas que os possibilite ler, ampliar e realizar releituras dessas práticas corporais, sejam elas de âmbito local, regional ou mundial.
Essa categoria articula -se às demais áreas do conhecimento pela própria essência da Cultura Corporal, constituída dialeticamente pelas relações humanas e materializada nas práticas corporais, experiências culturais, artísticas e literárias que integrem os campos de saberes e práticas por meio de projetos e/ou atividades de caráter inter e transdisciplinares.
Categoria de área : Jornalístico-Midiático Campo de saberes e práticas da Educação Física : Esta categoria requer atenção especial, devido à quantidade desenfreada de produção e circulação de textos e discursos presentes nas diversas mídias sobre os temas referentes à cultura corporal.
Portanto, é imprescindível pensar na recepção desses textos, preparar o aluno para discernir fatos, opiniões, fake news, isto é, promover uma postura crítica frente a estas temáticas.
Em relação às outras áreas, o diálogo direto com esta categoria pode se dar tanto pela abordagem de textos e discursos jornalísticos e midiáticos como pela análise/síntese dos acontecimentos e transformações sociais que permeiam as práticas corporais sob a influência da Indústria Cultural.
Categoria de área : Atuação na vida Pública Campo de saberes e práticas da Educação Física : Essa categoria propõe práticas pedagógicas que instrumentalizem e fomentem a proatividade dos estudantes de modo a intervir e transformar o contexto social, político e cultural em que os jovens estão inseridos, combatendo posicionamentos passivos e indiferentes frente aos problemas do cotidiano.
As atividades trabalhadas devem ter como objetivo desenvolver nos estudantes a compreensão do papel do poder público na garantia do acesso às práticas corporais, tais como as de esporte e lazer, de modo a atender seu protagonismo social.
Essa categoria possibilita também um trabalho interdisciplinar com atuação nos diferentes campos do conhecimento, portanto é imprescindível a abordagem dos objetos de conhecimentos da cultura corporal por meios diversificados que trabalhem a representatividade, a oratória, a expressividade, o protagonismo e a autonomia.
Um exemplo é o uso e abordagem de atividades ligadas às artes cênicas, workshops de dança ou expressão corporal, e atividades outras que possibilitem o seu empoderamento nas políticas públicas em esporte e lazer, bem como a temáticas relacionadas à saúde como direito inalienável ao exercício pleno de sua cidadania.
A partir dos campos de saberes e práticas do ensino, apresentados na área de Linguagens e suas Tecnologias, propõe -se a seguir um Organizador Curricular disposto por meio dos princípios curriculares norteadores da educação básica paraense, as categorias conceitual-metodológicas que fundamentam a área mencionada, além da relação entre as \textbf{competências} e habilidades específicas e os possíveis objetos de conhecimento que devem ser consolidados na formação geral básica.
Esse organizador pretende sintetizar a proposta curricular da área, bem como ser um dos pontos de partida para o planejamento integrado dos professores da área de Linguagens e suas tecnologias no Ensino Médio.

% matched lemmas: competência
\end{textsample}
